% =========================================
% CHAPTER 4: FRAMEWORK DESIGN
% =========================================
\chapter{Framework Design and Architecture}
\label{ch:framework}

This chapter describes the design of the testing framework. We run both rollup types against the same Layer~1 node and measure what actually happens under comparable conditions. The framework captures finality latency, recovery time after crashes, and throughput, and it can be reproduced with a single Docker command.

\section{Design Requirements}

Five requirements drove the design:

\textbf{Precise finality measurement.} Both rollups record when transactions are submitted and when they become irreversible. We want measured numbers from the running system, not theoretical estimates.

\textbf{Failure observation.} We simulate sequencer crashes and measure recovery time and data loss. A comparative framework that only tests the happy path misses the most architecturally interesting behavior.

\textbf{Shared Layer~1.} Both rollups run against the same Hardhat node. If they ran against separate nodes or at different times, observed differences could come from the environment rather than the architecture. Sharing a node eliminates that confound.

\textbf{Single-command deployment.} The full testbed starts with \texttt{docker-compose up -{}-build}. Reproducing the results should not require manual configuration steps that differ between machines.

\textbf{Sufficient instrumentation.} Every sequencer exports throughput, batch counts, finality latency, and—for ZK—proof generation time. These feed into Prometheus and Grafana for time-series analysis.

\section{System Architecture}

The framework has four components that work together to compare both rollup types.

The first component is a local Ethereum node running Hardhat~\cite{hardhat2024}. This node hosts both the ZKVerifier and OptimisticVerifier contracts. Because both contracts run on the same blockchain instance, any network delay or block timing variation affects both rollups equally. This prevents us from accidentally favoring one architecture.

The second component is two sequencers, each running as a separate Python process. They run within the same Docker~\cite{docker2024} container and connect to the shared Hardhat node. The ZK sequencer listens on port 5000; the Optimistic sequencer listens on port 5001. Running them as separate processes means if one crashes, the other continues—we can measure failures in isolation.

The third component is a web dashboard built with Flask~\cite{flask2024}. It shows real-time metrics from both sequencers and lets us trigger test scenarios. When we want to simulate a sequencer crash or inject a Layer~1 delay, we do it from the dashboard.

The fourth component is the metrics stack. Prometheus~\cite{prometheus2024} scrapes metrics from each sequencer (ports 9100 and 9101) and Grafana~\cite{grafana2024} visualizes the data. This gives us time-series records we can analyze later.

Figure~\ref{fig:architecture} shows how these pieces connect.

\begin{figure}[htbp]
    \centering
    \includegraphics[width=0.85\textwidth]{architecture/system-architecture.png}
    \caption{System architecture showing the four main components—Hardhat L1 node,
    ZK Sequencer (port 5000 / metrics 9100), Optimistic Sequencer (port 5001 / metrics 9101),
    Flask Dashboard (port 3000), Prometheus (port 9090), and Grafana (port 3001)—and
    how they connect inside the Docker network.}
    \label{fig:architecture}
\end{figure}

To deploy the framework, we run a single Docker Compose~\cite{compose2024} command. This starts all components at once. Table~\ref{tab:eval-dimensions} lists what we measure and how.

\begin{table}[h]
    \centering
    \begin{tabular}{>{\raggedright\arraybackslash}p{3.5cm}|>{\raggedright\arraybackslash}p{6.5cm}|l}
        \hline
        \textbf{What We Measure} & \textbf{Method} & \textbf{Unit} \\
        \hline
        Finality time & Record submission block and finality block & Blocks / seconds \\
        \hline
        Throughput & Count transactions per time window & Tx/second \\
        \hline
        Batching & Record batch size and timing & Tx per batch \\
        \hline
        Proof generation time & Time ZK sequencer spends computing commitment & Seconds (histogram) \\
        \hline
        L1 gas per submission & Gas from Hardhat receipts (local estimate only) & Gas units \\
        \hline
        State validity & Check if state roots match & Binary \\
        \hline
    \end{tabular}
    \caption{Measurement dimensions and methods}
    \label{tab:eval-dimensions}
\end{table}

\section{How Transactions Flow Through the Framework}

When we submit a transaction, here is what happens. The test scenario sends HTTP POST to one of the sequencers. The sequencer adds the transaction to a pending pool. When the pool reaches a threshold (100 transactions for ZK, 200 for Optimistic), the sequencer creates a batch. The sequencer computes the state root. For ZK, it also generates a proof. Then it sends the commitment to Layer~1.

The Layer~1 contract either verifies the proof (ZK) or opens a challenge window (Optimistic). For ZK, finality happens when the proof is verified and five additional blocks have passed. For Optimistic, finality happens when the challenge window closes and no dispute was filed.

We record timestamps at each step. The difference between submission block and finality block is the latency we measure. This latency includes both on-chain delays and off-chain computation.

\begin{figure}[htbp]
    \centering
    \includegraphics[width=0.70\textwidth]{architecture/batch-flow.png}
    \caption{The \texttt{\_batch()} function in both sequencers: the ZK path adds a mock proof hash and reloads \texttt{contracts.json} each call; the Optimistic path uses a preloaded handle and increments a batch counter.}
    \label{fig:batch-flow}
\end{figure}

\section{Why These Design Choices}

We use SHA-256 for proof generation rather than real zk-SNARKs. Real proofs require elliptic-curve pairing computations and, in practice, take seconds to minutes per batch on specialized hardware. SHA-256 lets us run many test scenarios quickly without that bottleneck. Both sequencers still generate a commitment and submit it to Layer~1—the architectural difference (proof required vs.\ not required) is preserved even though the proof itself is simplified. Production ZK~Rollups report proof generation times ranging from seconds to minutes per batch, so any overhead figures measured here are a lower bound.

The ZK batch size is 100 transactions and Optimistic is 200. Larger batches amortize the fixed cost of each Layer~1 submission across more transactions, which is good—but only when batch assembly itself is fast. With ZK, the commitment step is the bottleneck, so smaller batches keep per-transaction latency low. With Optimistic, there is no proof step, so the batching interval can be longer and the batch larger. These numbers reflect the trade-off each model makes.

The Optimistic challenge window is compressed to 10 blocks (~2 minutes on a Hardhat node) from the 7-day production value. Any evaluation framework needs to complete within a reasonable time. Ten blocks preserves the structural property—invalid state can be challenged before finality—while keeping test runs tractable.

Both sequencers run against the same Layer~1 node. Isolating each in its own node would introduce timing and block-production differences that could skew results. Shared execution means any observed latency or throughput difference is architectural, not environmental.

\section{Metrics We Export}

Each sequencer exports metrics to Prometheus on a dedicated port. The metrics are simple: transaction count, batch count, throughput, finality latency, and proof generation time (ZK only). We store these as histograms so we can compute percentiles—not just the average, but the P95 and P99 latencies, which matter more for user experience.

Table~\ref{tab:research-questions} shows how each research question maps to what we measure.

\begin{table}[h]
    \centering
    \begin{tabular}{l|>{\raggedright\arraybackslash}p{4.5cm}|>{\raggedright\arraybackslash\ttfamily\small}p{5cm}}
        \hline
        \textbf{Question} & \textbf{What We Measure} & \normalfont\textbf{Test File} \\
        \hline
        RQ1: Finality & Time from submission to finality &
            tests/integration/ \newline test\_dual\_sequencer.py \\
        \hline
        RQ2: Recovery & Time and cost to recover after crash &
            tests/adversarial/ \newline test\_sequencer\_crash.py \\
        \hline
        RQ3: Security cost & Proof generation time and gas overhead &
            tests/integration/ \newline test\_dual\_sequencer.py \\
        \hline
    \end{tabular}
    \caption{Research questions and corresponding test files}
    \label{tab:research-questions}
\end{table}

\subsection{Technology Stack}

Table~\ref{tab:tech-stack} lists the technologies used and why each was chosen. The criteria were simple: each tool had to be open-source, easy to run locally, and well enough documented that another researcher could replace it if needed.

\begin{table}[h]
    \centering
    \begin{tabular}{l|l|>{\raggedright\arraybackslash}p{5cm}}
        \hline
        \textbf{Component} & \textbf{Technology} & \textbf{Reason for choice} \\
        \hline
        Layer~1 node & Hardhat~\cite{hardhat2024} & Deterministic local blocks; no real ETH needed; fast reset between runs \\
        \hline
        Smart contracts & Solidity + ethers.js & Standard Ethereum toolchain; compatible with Hardhat out of the box \\
        \hline
        Sequencer processes & Python 3 + web3.py & Fast to iterate; logic is readable without compiled-language overhead \\
        \hline
        HTTP interface & Flask~\cite{flask2024} & Lightweight; transaction submission endpoint needs nothing heavier \\
        \hline
        Metrics export & Prometheus client & Each sequencer exposes counters and histograms on a dedicated port \\
        \hline
        Metrics dashboard & Grafana~\cite{grafana2024} & Standard time-series visualisation; queries written in PromQL \\
        \hline
        Deployment & Docker Compose~\cite{compose2024} & One command starts all services; environment variables control every parameter \\
        \hline
    \end{tabular}
    \caption{Technology stack and selection rationale}
    \label{tab:tech-stack}
\end{table}

\section{Configuration}

The framework parameters are configurable via environment variables. This allows researchers to test how different architectural choices affect performance.

\begin{table}[h]
    \centering
    \begin{tabular}{>{\ttfamily}l|l|>{\raggedright\arraybackslash}p{5.5cm}}
        \hline
        \normalfont\textbf{Parameter} & \textbf{Default} & \textbf{Purpose} \\
        \hline
        ZK\_BATCH\_SIZE & 100 & Transactions per ZK batch \\
        \hline
        ZK\_BATCH\_INTERVAL & 5 & Seconds between ZK batch checks \\
        \hline
        ZK\_GAS\_LIMIT & 200\,000 & Gas limit for ZK submissions (Hardhat estimate; not a mainnet figure) \\
        \hline
        OPT\_BATCH\_SIZE & 200 & Transactions per Optimistic batch \\
        \hline
        OPT\_BATCH\_INTERVAL & 8 & Seconds between Optimistic batch checks \\
        \hline
        OPT\_GAS\_LIMIT & 300\,000 & Gas limit for Optimistic submissions; raised from 150\,000 after measuring cold SSTORE cost of the 10-field Commitment struct \\
        \hline
        L1\_RPC & \normalfont\small http://localhost:8545 & Layer~1 RPC endpoint \\
        \hline
    \end{tabular}
    \caption{Configurable framework parameters}
    \label{tab:config-params}
\end{table}

These parameters can be set via Docker environment variables when starting the framework, allowing experimental variation without code changes. For example:

\begin{verbatim}
docker-compose up -d
# Now modify one parameter and re-run the full stack:
ZK_BATCH_SIZE=150 OPT_BATCH_INTERVAL=5 docker-compose up -d
\end{verbatim}

This approach lets researchers measure how batch size, polling frequency, and gas allocation affect finality latency, throughput, and cost. Changing any parameter and re-running the test suite provides immediate feedback on its impact.

\section{What This Chapter Establishes}

The framework gives us a controlled environment in which both rollup types run against identical conditions: the same Layer~1 node, the same hardware, the same network. The design choices documented here---batch sizes, challenge window compression, SHA-256 as a proof stand-in---are deliberate compromises between experimental practicality and architectural faithfulness. Each one is flagged where it appears so that results can be interpreted with appropriate scope.

Two intentional omissions are worth restating before moving to the code. First, transactions are not executed against a real state tree; state roots are derived by hashing batch data, not by applying each transaction to a storage trie. Second, ZK proofs are not verified cryptographically; the contract checks that a proof is present, not that it is correct. Both simplifications preserve the structural properties that matter for the comparison---ZK requires a proof at submission time, Optimistic requires a challenge window before finality---without the hardware and time cost that real cryptographic operations would impose.

Chapter~\ref{ch:implementation} describes how each component was built in code: what the contracts store, how the sequencers decide when to batch, and how the metrics layer records timestamps at each step.
